\documentclass[a4paper,10pt,brazil]{article}
\usepackage{provastyle}
\usepackage{menukeys}

\begin{document}
\makeexamtitlepage{UFOP}{BCC222: Programação Funcional}{José Romildo}{2026/1}{Primeira Prova}{04}
\begin{multicols}{2}
\questionTitle{1}{NoBuffering}
\begin{flushright}
  \footnotesize
  \textit{\textbf{8: Programas interativos}}\\
  \textit{c08-buffering-03}\\

\end{flushright}
Qual é o efeito de executar a ação abaixo no início de um programa?

\begin{minted}{haskell}
hSetBuffering stdout NoBuffering
\end{minted}

\begin{enumerate}
  \item Ela configura a saída padrão para não manter conteúdo em \emph{buffer}, fazendo com que as saídas sejam enviadas imediatamente.

  \item Ela impede que o programa escreva na saída padrão.

  \item Ela converte todo texto impresso para maiúsculas.

  \item Ela configura \haskellinline|stdout| para ler dados de um arquivo.

  \item Ela faz com que a entrada padrão seja lida caractere por caractere.


\end{enumerate}

\questionTitle{2}{Rastreamento do acumulador na leitura até zero}
\begin{flushright}
  \footnotesize
  \textit{\textbf{10: Ações de E/S recursivas}}\\
  \textit{c10-acumuladores-cauda-03}\\

\end{flushright}
Na função abaixo, o usuário digita a sequência \haskellinline|3, 4, 5, 0|, um valor por linha.

\begin{haskellcode}
lerESomar :: Integer -> IO Integer
lerESomar total = do
  n <- readLn
  if n == 0
    then return total
    else lerESomar (total + n)
\end{haskellcode}

Qual é o valor de \haskellinline|total| no momento em que o caso base é alcançado?

\begin{enumerate}
  \item 0

  \item 11

  \item 12

  \item 13

  \item 5


\end{enumerate}

\questionTitle{3}{Rastreamento da multiplicação recursiva}
\begin{flushright}
  \footnotesize
  \textit{\textbf{9: Funções recursivas}}\\
  \textit{c09-rastreamento-mul-01}\\

\end{flushright}
Considere:
\begin{haskellcode}
mul :: Int -> Int -> Int
mul m n
  | n == 0    = 0
  | n > 0     = m + mul m (n - 1)
  | otherwise = negate (mul m (negate n))
\end{haskellcode}
Qual é o valor de \haskellinline|mul -4 3|?

\begin{enumerate}
  \item \haskellinline|-16|

  \item \haskellinline|-12|

  \item \haskellinline|0|

  \item \haskellinline|12|

  \item \haskellinline|-7|


\end{enumerate}

\questionTitle{4}{Uso infixo de uma função prefixa}
\begin{flushright}
  \footnotesize
  \textit{\textbf{3: Expressões e definições}}\\
  \textit{c03-aplicacao-notacoes-04}\\

\end{flushright}
Qual alternativa é equivalente à expressão \haskellinline|div (sum [1,2,3,4]) 3| usando a função \haskellinline|div| em notação infixa?

\begin{enumerate}
  \item \haskellinline|sum [1,2,3,4] div 3|.

  \item \haskellinline|div `sum [1,2,3,4]` 3|.

  \item \haskellinline|3 `div` (sum [1,2,3,4])|.

  \item \haskellinline|(sum [1,2,3,4]) `div` 3|.

  \item \haskellinline|(sum [1,2,3,4], 3) `div`|.


\end{enumerate}

\questionTitle{5}{Modelagem de biblioteca}
\begin{flushright}
  \footnotesize
  \textit{\textbf{6: Estruturas de dados básicas}}\\
  \textit{c06-modelagem-01}\\

\end{flushright}
Considere as definições abaixo:
\begin{minted}{haskell}
type Titulo = String
type Autor  = String
type Ano    = Maybe Int
type Livro  = (Titulo, Autor, Ano)
\end{minted}
Qual definição representa corretamente uma biblioteca como coleção de livros?

\begin{enumerate}
  \item \haskellinline{type Biblioteca = Livro}

  \item \haskellinline{type Biblioteca = [Livro]}

  \item \haskellinline{type Biblioteca = Maybe Livro}

  \item \haskellinline{type Biblioteca = (Livro, Livro)}

  \item \haskellinline{type Biblioteca = [Titulo]}


\end{enumerate}

\questionTitle{6}{Ausência de segurança estrita em sinônimos}
\begin{flushright}
  \footnotesize
  \textit{\textbf{4: Tipos de Dados}}\\
  \textit{c04-sinonimos-equivalencia-01}\\

\end{flushright}
Considere:
\begin{haskellcode}
type Nome = String
type Endereco = String
geraEtiqueta :: Nome -> Endereco -> String
geraEtiqueta n e = "Para: " ++ n ++ " - Env: " ++ e
\end{haskellcode}
Se o programador chamar \haskellinline{geraEtiqueta "Rua X" "João"} (invertendo os argumentos), o que acontece?

\begin{enumerate}
  \item O compilador rejeita o código com um erro de tipo, protegendo contra a inversão.

  \item O código compila sem erros ou avisos, pois \haskellinline{Nome} e \haskellinline{Endereco} são o mesmo tipo que \haskellinline{String}.

  \item O compilador emite um aviso, mas gera o executável.

  \item O código compila, mas lança uma exceção em tempo de execução.

  \item Ocorre um erro sintático porque sinônimos exigem construtores (ex: \haskellinline{Nome "João"}).


\end{enumerate}

\questionTitle{7}{Sintaxe correta com guardas}
\begin{flushright}
  \footnotesize
  \textit{\textbf{5: Expressão condicional}}\\
  \textit{c05-guardas-sintaxe-semantica-01}\\

\end{flushright}
Qual das definições abaixo está escrita corretamente usando guardas em Haskell?

\begin{enumerate}
  \item \begin{haskellcode}
abs' n
  n < 0 -> -n
  otherwise -> n
\end{haskellcode}

  \item \begin{haskellcode}
abs' n =
  | n < 0     = -n
  | otherwise = n
\end{haskellcode}

  \item \begin{haskellcode}
abs' n
  | n < 0     = -n
  | otherwise = n
\end{haskellcode}

  \item \begin{haskellcode}
abs' n
  if n < 0 = -n
  else     = n
\end{haskellcode}

  \item \begin{haskellcode}
abs' n
  | n < 0 then -n
  | otherwise then n
\end{haskellcode}


\end{enumerate}

\questionTitle{8}{Fundação teórica: Cálculo Lambda}
\begin{flushright}
  \footnotesize
  \textit{\textbf{1: Paradigmas de programação}}\\
  \textit{c01-historia-evolucao-01}\\

\end{flushright}
Na década de 1930, antes mesmo da existência dos computadores eletrônicos construídos em silício, foi desenvolvido um sistema formal matemático que provou ser tão poderoso e universal quanto a Máquina de Turing. Este sistema lançou a fundação teórica de toda a programação funcional moderna.

Qual é o nome deste sistema matemático e quem foi o seu criador?

\begin{enumerate}
  \item A Arquitetura Pura, desenvolvida por John von Neumann.

  \item A Teoria da Avaliação Preguiçosa (\emph{Lazy Evaluation}), formalizada por David Turner.

  \item A Álgebra Booleana, desenvolvida por George Boole.

  \item O Sistema de Inferência de Tipos de Primeira Ordem, desenvolvido por Robin Milner.

  \item O Cálculo Lambda, desenvolvido por Alonzo Church.


\end{enumerate}

\questionTitle{9}{Rastreamento de recursividade mútua}
\begin{flushright}
  \footnotesize
  \textit{\textbf{9: Funções recursivas}}\\
  \textit{c09-formas-mutua-01}\\

\end{flushright}
Considere:
\begin{haskellcode}
par n | n == 0 = True
      | otherwise = impar (n - 1)

impar n | n == 0 = False
        | otherwise = par (n - 1)
\end{haskellcode}
Qual é o valor de \haskellinline|impar 4|?

\begin{enumerate}
  \item \haskellinline|False|

  \item erro de tipo

  \item não termina

  \item \haskellinline|4|

  \item \haskellinline|True|


\end{enumerate}

\questionTitle{10}{Avaliação básica de guardas}
\begin{flushright}
  \footnotesize
  \textit{\textbf{5: Expressão condicional}}\\
  \textit{c05-guardas-avaliacao-01}\\

\end{flushright}
Considere a função:
\begin{haskellcode}
f x
  | x > 10 = 'A'
  | x > 5  = 'B'
  | otherwise = 'C'
\end{haskellcode}
Qual é o valor de \haskellinline|f 7|?

\begin{enumerate}
  \item Ocorre um erro de guardas não exaustivas.

  \item \haskellinline|'A'|

  \item \haskellinline|'B'|

  \item \haskellinline|'C'|

  \item \haskellinline|"B"|


\end{enumerate}

\questionTitle{11}{Ações de saída padrão}
\begin{flushright}
  \footnotesize
  \textit{\textbf{8: Programas interativos}}\\
  \textit{c08-entrada-saida-03}\\

\end{flushright}
Considere a função de saída \haskellinline|putChar|. Qual alternativa descreve corretamente seu tipo e seu comportamento?

\begin{enumerate}
  \item \haskellinline|String -> IO ()| — imprime uma string inteira como se fosse um caractere.

  \item \haskellinline|String -> String| — transforma texto sem construir nenhuma ação de E/S.

  \item \haskellinline|Char -> IO ()| — constrói uma ação que imprime um caractere.

  \item \haskellinline|IO Char| — lê um caractere da entrada padrão.

  \item \haskellinline|IO String| — lê da entrada padrão o texto que deveria ser impresso.


\end{enumerate}

\questionTitle{12}{Tupla, lista e agrupamento}
\begin{flushright}
  \footnotesize
  \textit{\textbf{6: Estruturas de dados básicas}}\\
  \textit{c06-tuplas-sintaxe-01}\\

\end{flushright}
Qual das expressões abaixo representa uma tupla de exatamente dois elementos?

\begin{enumerate}
  \item \haskellinline{(42, "ok")}

  \item \haskellinline{("ok")}

  \item \haskellinline{42, "ok"}

  \item \haskellinline{(42)}

  \item \haskellinline{[42, "ok"]}


\end{enumerate}

\questionTitle{13}{map show aplicado a literais numéricos}
\begin{flushright}
  \footnotesize
  \textit{\textbf{7: Polimorfismo}}\\
  \textit{cap07-expressao-map-show-numerico}\\

\end{flushright}
Considere a expressão:

\begin{minted}{haskell}
map show [1, 2, 3]
\end{minted}

Qual análise é a mais precisa?

\begin{enumerate}
  \item O resultado tem tipo \haskellinline{[String]}; antes da resolução padrão, os elementos da lista são restringidos por \haskellinline{Num a} e \haskellinline{Show a}.
  \item O resultado tem tipo \haskellinline{String}, pois \haskellinline{map} concatena automaticamente a lista resultante.
  \item O resultado tem tipo \haskellinline{[Int]}, pois \haskellinline{show} preserva o tipo do argumento.
  \item A expressão tem tipo \haskellinline{Show a => [a]}, pois \haskellinline{map} não altera o tipo da lista.
  \item A expressão é impossível, pois \haskellinline{show} não pode ser passada como argumento para \haskellinline{map}.

\end{enumerate}

\questionTitle{14}{Comando :type no GHCi}
\begin{flushright}
  \footnotesize
  \textit{\textbf{4: Tipos de Dados}}\\
  \textit{c04-ghci-consulta-01}\\

\end{flushright}
Qual comando do GHCi exibe o tipo de uma expressão sem avaliá-la?

\begin{enumerate}
  \item \texttt{:kind}

  \item \texttt{:info}

  \item \texttt{:show}

  \item \texttt{:load}

  \item \texttt{:type}


\end{enumerate}

\questionTitle{15}{Saída produzida por mostraLista}
\begin{flushright}
  \footnotesize
  \textit{\textbf{10: Ações de E/S recursivas}}\\
  \textit{c10-rastreamento-03}\\

\end{flushright}
Considere:

\begin{haskellcode}
mostraLista :: Show a => [a] -> IO ()
mostraLista lista
  | null lista = return ()
  | otherwise = do
      print (head lista)
      mostraLista (tail lista)
\end{haskellcode}

Qual é a saída de \haskellinline|mostraLista [2,4,6]|?

\begin{enumerate}
  \item Apenas \haskellinline|2| é impresso, pois a chamada recursiva é ignorada.

  \item \haskellinline|[2,4,6]| em uma única linha.

  \item Nada é impresso, pois a função retorna \haskellinline|()|.

  \item Os valores \haskellinline|6|, \haskellinline|4| e \haskellinline|2| são impressos, um por linha.

  \item Os valores \haskellinline|2|, \haskellinline|4| e \haskellinline|6| são impressos, um por linha, nessa ordem.


\end{enumerate}

\questionTitle{16}{Definição local em programa interativo}
\begin{flushright}
  \footnotesize
  \textit{\textbf{8: Programas interativos}}\\
  \textit{c08-programas-completos-03}\\

\end{flushright}
Analise:

\begin{minted}{haskell}
main :: IO ()
main = do
  n <- readLn :: IO Int
  let dobro = 2 * n
  print dobro
\end{minted}

Qual afirmação descreve corretamente a linha com \haskellinline|let|?

\begin{enumerate}
  \item Ela executa uma nova ação de E/S para ler o valor de \haskellinline|dobro|.

  \item Ela imprime automaticamente o dobro antes da linha \haskellinline|print dobro|.

  \item Ela só seria válida se fosse escrita como \haskellinline|dobro <- 2 * n|.

  \item Ela transforma \haskellinline|dobro| em uma ação do tipo \haskellinline|IO Int|.

  \item Ela define localmente o valor puro \haskellinline|dobro|, usando o valor \haskellinline|n| já obtido pela ação \haskellinline|readLn|.


\end{enumerate}

\questionTitle{17}{Comparação conceitual entre let e where}
\begin{flushright}
  \footnotesize
  \textit{\textbf{3: Expressões e definições}}\\
  \textit{c03-let-escopo-04}\\

\end{flushright}
Qual alternativa apresenta uma diferença conceitual correta entre \haskellinline|let| e \haskellinline|where|?

\begin{enumerate}
  \item Não há diferença: \haskellinline|let| e \haskellinline|where| são apenas grafias alternativas para a mesma construção sintática.

  \item \haskellinline|where| é uma expressão, enquanto \haskellinline|let| só pode aparecer no fim de uma função.

  \item \haskellinline|let| permite apenas uma definição local, enquanto \haskellinline|where| permite apenas literais numéricos.

  \item \haskellinline|let ... in ...| é uma expressão e pode aparecer onde expressões são permitidas; \haskellinline|where| é uma cláusula anexada a uma equação.

  \item \haskellinline|where| sempre tem escopo global, e \haskellinline|let| sempre tem escopo de módulo.


\end{enumerate}

\questionTitle{18}{Construção e execução de ações}
\begin{flushright}
  \footnotesize
  \textit{\textbf{8: Programas interativos}}\\
  \textit{c08-modelo-io-03}\\

\end{flushright}
Considere a definição:

\begin{minted}{haskell}
saudacao :: IO ()
saudacao = putStrLn "Olá!"
\end{minted}

O que essa definição faz no momento em que o módulo é carregado?

\begin{enumerate}
  \item Ela define uma função pura do tipo \haskellinline|String -> String| que transforma a mensagem antes de imprimi-la.

  \item Ela executa a ação uma vez e guarda o resultado textual \haskellinline|"Olá!"| para reutilização futura.

  \item Ela associa o nome \haskellinline|saudacao| a uma ação de E/S. A mensagem só será impressa quando essa ação for executada.

  \item Ela imprime imediatamente \haskellinline|"Olá!"|, pois toda definição de tipo \haskellinline|IO ()| é executada ao carregar o módulo.

  \item Ela cria uma exceção, pois \haskellinline|putStrLn| só pode aparecer diretamente dentro de \haskellinline|main|.


\end{enumerate}

\questionTitle{19}{Leitura de condicional aninhada}
\begin{flushright}
  \footnotesize
  \textit{\textbf{5: Expressão condicional}}\\
  \textit{c05-if-aninhamento-01}\\

\end{flushright}
Considere a definição:
\begin{haskellcode}
sinal n = if n < 0 then -1 else if n > 0 then 1 else 0
\end{haskellcode}
Qual é o resultado de \haskellinline|sinal 0|?

\begin{enumerate}
  \item Ocorre um erro, pois o segundo \haskellinline|if| não tem \haskellinline|else|.

  \item \haskellinline|False|

  \item \haskellinline|0|

  \item \haskellinline|-1|

  \item \haskellinline|1|


\end{enumerate}

\questionTitle{20}{Identificação de forma normal}
\begin{flushright}
  \footnotesize
  \textit{\textbf{3: Expressões e definições}}\\
  \textit{c03-layout-avaliacao-04}\\

\end{flushright}
Qual das alternativas abaixo está em forma normal no sentido apresentado no capítulo, isto é, não precisa de mais reduções para ser um valor?

\begin{enumerate}
  \item \haskellinline|sqrt 16|.

  \item \haskellinline|(\x -> x) 5|.

  \item \haskellinline|let x = 3 in x + 1|.

  \item \haskellinline|1 + 2 * 3|.

  \item \haskellinline|7|.


\end{enumerate}

\questionTitle{21}{Definição formal de transparência referencial}
\begin{flushright}
  \footnotesize
  \textit{\textbf{1: Paradigmas de programação}}\\
  \textit{c01-transparencia-referencial-01}\\

\end{flushright}
Na teoria de linguagens de programação funcionais, o que significa afirmar que uma dada expressão possui \textbf{transparência referencial}?

\begin{enumerate}
  \item Significa que o código-fonte da função não pode ser ofuscado pelo compilador, permanecendo \emph{transparente} e legível para fins de auditoria no binário.

  \item Significa que o sistema de inferência de tipos consegue deduzir a assinatura da função mesmo que o programador omita essa informação no cabeçalho.

  \item Significa que a expressão pode ser substituída diretamente pelo seu valor computado correspondente em qualquer ponto do programa, sem alterar o comportamento global do sistema.

  \item Significa que a função tem permissão explícita para ler variáveis do escopo global mutável, desde que não as modifique durante a execução.

  \item Significa que a função permite que referências de ponteiros sejam passadas como argumentos em vez de cópias de valores, economizando alocação de memória.


\end{enumerate}

\questionTitle{22}{Consulta com zip e lookup}
\begin{flushright}
  \footnotesize
  \textit{\textbf{6: Estruturas de dados básicas}}\\
  \textit{c06-lookup-associacao-01}\\

\end{flushright}
Considere a expressão abaixo.
\begin{minted}{haskell}
let chaves = [5,6,7]
    valores = ["p","q","r"]
    tabela = zip chaves valores
in lookup 5 tabela
\end{minted}
Qual é o resultado final?

\begin{enumerate}
  \item Ocorre um erro de tipo.

  \item \haskellinline{ Just "q" }

  \item \haskellinline{ Just "p" }

  \item \haskellinline{ "p" }

  \item \haskellinline{ Nothing }


\end{enumerate}

\questionTitle{23}{Inferência de tipo de not True}
\begin{flushright}
  \footnotesize
  \textit{\textbf{4: Tipos de Dados}}\\
  \textit{c04-inferencia-basica-01}\\

\end{flushright}
Qual é o tipo inferido para a expressão \haskellinline{not True}?

\begin{enumerate}
  \item \haskellinline{Bool}

  \item \haskellinline{True}

  \item \haskellinline{False}

  \item \haskellinline{Boolean}

  \item \haskellinline{Bool -> Bool}


\end{enumerate}

\questionTitle{24}{Classes de tipo e interfaces de POO}
\begin{flushright}
  \footnotesize
  \textit{\textbf{7: Polimorfismo}}\\
  \textit{cap07-comparacao-typeclass-poo}\\

\end{flushright}
Uma analogia comum aproxima classes de tipo em Haskell de interfaces em linguagens orientadas a objetos. Qual afirmação evita a simplificação excessiva?
\begin{enumerate}
  \item Classes de tipo não têm qualquer relação com operações; elas servem apenas para documentação.
  \item Interfaces de POO e classes de tipo são idênticas porque ambas exigem que todos os métodos sejam chamados dinamicamente.
  \item A analogia é útil porque ambas especificam operações disponíveis, mas classes de tipo são resolvidas por instâncias e restrições de tipos, não por subtipagem nominal de objetos mutáveis.
  \item Classes de tipo só existem para simular herança múltipla em Haskell.
  \item Classes de tipo são exatamente classes de POO, com campos, construtores e herança de implementação.

\end{enumerate}

\questionTitle{25}{Tipo e propósito de fromIntegral}
\begin{flushright}
  \footnotesize
  \textit{\textbf{7: Polimorfismo}}\\
  \textit{cap07-conversao-fromintegral-alcance}\\

\end{flushright}
Qual alternativa descreve corretamente a função \haskellinline{fromIntegral}?
\begin{enumerate}
  \item Ela converte de um tipo \haskellinline{Integral} para qualquer tipo de destino que seja instância de \haskellinline{Num}.
  \item Ela realiza coerção implícita automática em todos os operadores numéricos.
  \item Ela converte qualquer \haskellinline{Fractional} diretamente para \haskellinline{Int}, arredondando para o inteiro mais próximo.
  \item Ela converte strings numéricas para inteiros, substituindo \haskellinline{read}.
  \item Ela é usada apenas para converter \haskellinline{Double} em \haskellinline{Float}.

\end{enumerate}

\questionTitle{26}{Leitura de quantidade fixa com n não positivo}
\begin{flushright}
  \footnotesize
  \textit{\textbf{10: Ações de E/S recursivas}}\\
  \textit{c10-lacos-io-03}\\

\end{flushright}
Analise a função:

\begin{haskellcode}
leLista :: Integer -> IO [Integer]
leLista n
  | n <= 0    = return []
  | otherwise = do
      x <- readLn
      xs <- leLista (n - 1)
      return (x:xs)
\end{haskellcode}

Se \haskellinline|leLista (-2)| for executada, o que acontece?

\begin{enumerate}
  \item Nenhum valor é lido; a ação produz a lista vazia por meio de \haskellinline|return []|.

  \item A função tenta ler dois valores e depois devolve uma lista vazia.

  \item O programa falha por erro de tipo, pois \haskellinline|Integer| não pode ser comparado com zero.

  \item A função entra em recursão infinita, pois \haskellinline|n| nunca se torna igual a zero.

  \item A função devolve \haskellinline|[-2]|, pois o valor inicial de \haskellinline|n| é incluído na lista.


\end{enumerate}

\questionTitle{27}{Superclasses relevantes de RealFrac}
\begin{flushright}
  \footnotesize
  \textit{\textbf{7: Polimorfismo}}\\
  \textit{cap07-hierarquia-realfrac}\\

\end{flushright}
A classe \haskellinline{RealFrac} aparece em funções como \haskellinline{floor}, \haskellinline{ceiling} e \haskellinline{round}. Qual descrição da hierarquia é correta?
\begin{enumerate}
  \item \haskellinline{RealFrac} exige, entre outras relações, que o tipo seja compatível tanto com \haskellinline{Real} quanto com \haskellinline{Fractional}.
  \item \haskellinline{RealFrac} é usada apenas para inteiros, por isso \haskellinline{Double} não pertence a essa classe.
  \item \haskellinline{RealFrac} elimina a necessidade de \haskellinline{Ord}, pois arredondamento não envolve ordem.
  \item \haskellinline{RealFrac} é uma classe de POO usada para armazenar campos mutáveis.
  \item \haskellinline{RealFrac} é superclasse de \haskellinline{Num}, portanto todo tipo numérico é automaticamente \haskellinline{RealFrac}.

\end{enumerate}

\questionTitle{28}{Ambiguidade de show [] e EDR}
\begin{flushright}
  \footnotesize
  \textit{\textbf{7: Polimorfismo}}\\
  \textit{cap07-defaulting-show-lista-vazia}\\

\end{flushright}
Considere a definição:

\begin{minted}{haskell}
w = show []
\end{minted}

Qual comparação entre compilação normal e GHCi com \haskellinline{ExtendedDefaultRules} ligada é correta?

\begin{enumerate}
  \item O Relatório escolhe \haskellinline{Integer} para o elemento da lista, pois \haskellinline{Integer} é o primeiro default.
  \item A expressão só compila se a lista vazia for convertida com \haskellinline{fromInteger}.
  \item A expressão sempre é rejeitada, inclusive no GHCi com EDR, pois listas vazias não têm tipo.
  \item Pela regra do Relatório, a expressão é ambígua por não haver classe numérica; no GHCi com EDR, ela pode ser resolvida escolhendo \haskellinline{a ~ ()}.
  \item A expressão sempre tem tipo \haskellinline{[String]}, pois \haskellinline{show} retorna listas.

\end{enumerate}

\questionTitle{29}{Significado de otherwise}
\begin{flushright}
  \footnotesize
  \textit{\textbf{5: Expressão condicional}}\\
  \textit{c05-otherwise-01}\\

\end{flushright}
O que \haskellinline|otherwise| representa em uma definição com guardas?

\begin{enumerate}
  \item Uma palavra reservada que faz a avaliação voltar para a equação anterior.

  \item Um nome reservado equivalente a \haskellinline|False|.

  \item Um nome predefinido equivalente a \haskellinline|True|, normalmente usado como caso padrão.

  \item Uma variável criada automaticamente com o valor da guarda anterior.

  \item Uma forma especial de indicar erro em tempo de execução.


\end{enumerate}

\questionTitle{30}{Tipagem estática vs. dinâmica}
\begin{flushright}
  \footnotesize
  \textit{\textbf{4: Tipos de Dados}}\\
  \textit{c04-checagem-estatica-dinamica-01}\\

\end{flushright}
Qual é a principal diferença entre um sistema de tipos estático (como Haskell) e um dinâmico (como Python)?

\begin{enumerate}
  \item Na tipagem estática, os tipos podem mudar durante a execução; na dinâmica, são fixos.

  \item A tipagem estática é mais lenta, mas mais flexível que a dinâmica.

  \item Na tipagem estática, a verificação de tipos ocorre em tempo de compilação; na dinâmica, em tempo de execução.

  \item Não há diferença fundamental; são apenas termos diferentes para o mesmo conceito.

  \item Apenas linguagens compiladas podem ter tipagem estática.


\end{enumerate}

\questionTitle{31}{Identificadores simbólicos}
\begin{flushright}
  \footnotesize
  \textit{\textbf{3: Expressões e definições}}\\
  \textit{c03-arquivos-identificadores-04}\\

\end{flushright}
Qual alternativa apresenta um identificador simbólico apropriado para definir um novo operador infixo comum, sem usar a forma reservada a construtores?

\begin{enumerate}
  \item \haskellinline|A++|.

  \item \haskellinline|:++|.

  \item \haskellinline|<++>|.

  \item \haskellinline|op+|.

  \item \haskellinline|let|.


\end{enumerate}

\questionTitle{32}{Rastreamento de acumulador em potência de dois}
\begin{flushright}
  \footnotesize
  \textit{\textbf{9: Funções recursivas}}\\
  \textit{c09-cauda-pdois-01}\\

\end{flushright}
Considere a função auxiliar:
\begin{haskellcode}
pdois' x y
  | x == 0 = y
  | x > 0  = pdois' (x - 1) (2 * y)
\end{haskellcode}
A chamada \haskellinline|pdois' 4 1| passa imediatamente para qual chamada?

\begin{enumerate}
  \item \haskellinline|pdois' 3 2|

  \item \haskellinline|pdois' 16 1|

  \item \haskellinline|pdois' 4 2|

  \item \haskellinline|pdois' 3 1|

  \item \haskellinline|pdois' 2 3|


\end{enumerate}

\questionTitle{33}{Detalhes profundos com o comando :info}
\begin{flushright}
  \footnotesize
  \textit{\textbf{2: Ambiente de desenvolvimento}}\\
  \textit{c02-ghci-inspecao-04}\\

\end{flushright}
Considere que o aluno já sabe o tipo do operador de adição (\haskellinline{+}), mas deseja investigar em qual módulo interno ele foi definido e qual é o nível matemático de sua precedência (\emph{fixity}). Qual comando do GHCi fornece essa camada extra de metadados?

\begin{enumerate}
  \item O comando \texttt{:show} (ou \texttt{:s}).

  \item O comando \texttt{:browse} (ou \texttt{:b}).

  \item O comando \texttt{:help} (ou \texttt{:?}).

  \item O comando \texttt{:type} (ou \texttt{:t}).

  \item O comando \texttt{:info} (ou \texttt{:i}).


\end{enumerate}

\questionTitle{34}{Visibilidade em guardas e resultados}
\begin{flushright}
  \footnotesize
  \textit{\textbf{5: Expressão condicional}}\\
  \textit{c05-where-escopo-01}\\

\end{flushright}
Em uma equação com guardas e uma cláusula \haskellinline|where|, onde são visíveis as definições locais?

\begin{enumerate}
  \item Apenas na primeira guarda da equação.

  \item Nas guardas e nas expressões de resultado da mesma equação.

  \item Apenas nas expressões de resultado, nunca nas guardas.

  \item Apenas nas guardas, nunca nas expressões de resultado.

  \item Em toda a função, inclusive em outras equações com o mesmo nome.


\end{enumerate}

\questionTitle{35}{Diagnóstico: Erro de diretório no load}
\begin{flushright}
  \footnotesize
  \textit{\textbf{2: Ambiente de desenvolvimento}}\\
  \textit{c02-fluxo-trabalho-04}\\

\end{flushright}
Um aluno salva o seu código sob o nome \texttt{Pratica.hs} na pasta \emph{Documentos} e, na sequência, abre um terminal recém-iniciado e digita o comando \texttt{:load Pratica.hs} no GHCi. O interpretador retorna o erro de sistema \emph{Target "Pratica.hs" is not a module name or a source file}. 

Qual é o passo crucial que foi omitido no fluxo de inicialização do ambiente?

\begin{enumerate}
  \item O GHCi obriga que o arquivo possua a extensão binária \texttt{.bin} ao invés da extensão de texto plano \texttt{.hs}.

  \item O aluno deveria ter forçado a importação global com o comando de sistema \texttt{:sudo load Pratica.hs}.

  \item O aluno não possuía a licença corporativa do \emph{Haskell Language Server} ativada para permitir a leitura remota de discos.

  \item O aluno nomeou o arquivo com a primeira letra em maiúscula (\texttt{P}), o que é estritamente proibido pela especificação sintática de módulos.

  \item O aluno esqueceu de garantir que o terminal do sistema operacional estava no mesmo diretório (pasta) onde o arquivo foi salvo antes de chamar o executável \texttt{ghci}.


\end{enumerate}

\questionTitle{36}{Remoção de módulos do escopo}
\begin{flushright}
  \footnotesize
  \textit{\textbf{2: Ambiente de desenvolvimento}}\\
  \textit{c02-ghci-modulos-04}\\

\end{flushright}
O atalho \texttt{:m} (abreviação de \texttt{:module}) do GHCi oferece uma flexibilidade adicional na gestão do escopo que a palavra-chave tradicional \haskellinline{import} não possui no terminal. Qual é essa flexibilidade?

\begin{enumerate}
  \item O comando \texttt{:m} permite a utilização do prefixo de subtração (\texttt{-}), como em \texttt{:m -Data.List}, propiciando a remoção rápida de módulos do escopo interativo atual.

  \item O comando \texttt{:m} contorna a necessidade de conectividade com a internet ao buscar módulos indisponíveis no cache.

  \item O comando \texttt{:m} força a avaliação estrita de toda a biblioteca sem utilizar as otimizações preguiçosas.

  \item O comando \texttt{:m} executa as funções do módulo importado em \emph{threads} separadas visando alto paralelismo.

  \item O comando \texttt{:m} realiza um \emph{dump} da memória do GHCi e grava o ambiente ativo em um novo arquivo \texttt{.hs}.


\end{enumerate}

\questionTitle{37}{Duas formas de construir uma lista lida}
\begin{flushright}
  \footnotesize
  \textit{\textbf{10: Ações de E/S recursivas}}\\
  \textit{c10-listas-ordem-04}\\

\end{flushright}
Compare as duas versões:

\begin{haskellcode}
-- Versão A
lerA :: IO [Integer]
lerA = do
  x <- readLn
  if x == 0
    then return []
    else do
      xs <- lerA
      return (x:xs)

-- Versão B
lerB :: [Integer] -> IO [Integer]
lerB xs = do
  x <- readLn
  if x == 0
    then return (reverse xs)
    else lerB (x:xs)
\end{haskellcode}

Qual alternativa descreve corretamente a diferença entre elas?

\begin{enumerate}
  \item A versão A inclui o sentinela na lista final, enquanto a versão B o descarta.

  \item As duas versões são idênticas quanto à ordem de construção, pois ambas usam o operador \haskellinline|:| no mesmo ponto.

  \item A versão B não precisa de caso base porque o acumulador substitui a condição de parada.

  \item A versão A devolve uma lista em ordem inversa e a versão B devolve sempre a lista vazia.

  \item A versão A reconstrói a lista ao retornar das chamadas recursivas; a versão B acumula na descida e inverte a lista ao final.


\end{enumerate}

\questionTitle{38}{Associatividade da seta de tipo}
\begin{flushright}
  \footnotesize
  \textit{\textbf{4: Tipos de Dados}}\\
  \textit{c04-funcao-associatividade-01}\\

\end{flushright}
Devido à associatividade à direita da seta de tipo (\haskellinline{->}), como a assinatura \haskellinline{Int -> Bool -> String} é implicitamente parentizada?

\begin{enumerate}
  \item \haskellinline{(Int -> Bool) -> String}

  \item \haskellinline{Int -> Bool -> (String)}

  \item A seta não possui associatividade definida.

  \item \haskellinline{(Int, Bool) -> String}

  \item \haskellinline{Int -> (Bool -> String)}


\end{enumerate}

\questionTitle{39}{Projeto de potenciação recursiva}
\begin{flushright}
  \footnotesize
  \textit{\textbf{9: Funções recursivas}}\\
  \textit{c09-projeto-potencia-01}\\

\end{flushright}
Qual definição calcula \haskellinline|x| elevado a um expoente natural \haskellinline|n| por multiplicações sucessivas?

\begin{enumerate}
  \item \begin{haskellcode}
potencia x n
  | n == 0 = 1
  | n > 0  = x + potencia x (n - 1)
\end{haskellcode}

  \item \begin{haskellcode}
potencia x n
  | n == 0 = 0
  | n > 0  = x * potencia x (n - 1)
\end{haskellcode}

  \item \begin{haskellcode}
potencia x n
  | n == 0 = 1
  | n > 0  = x * potencia x (n - 1)
\end{haskellcode}

  \item \begin{haskellcode}
potencia x n
  | n == 0 = 1
  | n > 0  = potencia x (n + 1)
\end{haskellcode}

  \item \begin{haskellcode}
potencia x n
  | n == 0 = x
  | n > 0  = x * potencia x n
\end{haskellcode}


\end{enumerate}

\questionTitle{40}{Mecânica da execução: Atribuição vs. Avaliação}
\begin{flushright}
  \footnotesize
  \textit{\textbf{1: Paradigmas de programação}}\\
  \textit{c01-mecanismo-reducao-01}\\

\end{flushright}
No paradigma imperativo (como na linguagem C), a computação avança predominantemente através da \textbf{execução de comandos} que alteram o estado da memória. Em contraste, qual é o mecanismo nuclear que impulsiona a computação em um programa funcional puro como Haskell?

\begin{enumerate}
  \item A validação de axiomas através de um motor de inferência lógica operando sobre um banco de dados de regras.

  \item A modificação sequencial de valores armazenados em registradores de variáveis estáticas no escopo global do módulo compilado.

  \item O despacho dinâmico de métodos orientados a eventos, desencadeados por interrupções do sistema operacional.

  \item A mutação explícita de estruturas de dados iteráveis através da alocação manual de ponteiros em tempo de execução.

  \item A avaliação e redução de expressões através da aplicação de funções, substituindo parâmetros até que se atinja um valor normal (final).


\end{enumerate}

\questionTitle{41}{Média de sequência vazia}
\begin{flushright}
  \footnotesize
  \textit{\textbf{10: Ações de E/S recursivas}}\\
  \textit{c10-bordas-robustez-03}\\

\end{flushright}
Considere o padrão usado no exercício da média:

\begin{haskellcode}
main :: IO ()
main = do
  lista <- lerLista
  print (sum lista / fromIntegral (length lista))

lerLista :: IO [Double]
lerLista = do
  x <- readLn
  if x < 0
    then return []
    else do
      xs <- lerLista
      return (x:xs)
\end{haskellcode}

Qual caso de borda precisa ser tratado explicitamente?

\begin{enumerate}
  \item A presença de zero na lista, pois zero é sempre interpretado como sentinela nessa função.

  \item A lista com um único valor, pois \haskellinline|sum| não funciona com listas unitárias.

  \item A lista em ordem crescente, pois \haskellinline|sum| exige que a entrada esteja ordenada.

  \item Qualquer lista com valores positivos, pois \haskellinline|fromIntegral| só aceita números negativos.

  \item A sequência vazia, pois se o primeiro valor digitado for negativo, \haskellinline|length lista| será zero e a média não estará bem definida.


\end{enumerate}

\questionTitle{42}{Inferência com sqrt, guardas e comparação}
\begin{flushright}
  \footnotesize
  \textit{\textbf{7: Polimorfismo}}\\
  \textit{cap07-inferencia-raizes-equacao}\\

\end{flushright}
Considere a definição:

\begin{minted}{haskell}
raizes a b c
  | delta > 0  = [(-b + sqrt delta)/(2*a)
                 ,(-b - sqrt delta)/(2*a)]
  | delta == 0 = [-b/(2*a)]
  | otherwise  = []
  where
    delta = b^2 - 4*a*c
\end{minted}

Qual é a assinatura de tipo mais geral?

\begin{enumerate}
  \item \mintinline{haskell}{raizes :: Integral a => a -> a -> a -> [a]}
  \item \mintinline{haskell}{raizes :: (Ord a, Floating a) => a -> a -> a -> [a]}
  \item \mintinline{haskell}{raizes :: Floating a => a -> a -> a -> Bool}
  \item \mintinline{haskell}{raizes :: Double -> Double -> Double -> Bool}
  \item \mintinline{haskell}{raizes :: (Ord a, Num a) => a -> a -> a -> [a]}

\end{enumerate}

\questionTitle{43}{Definição de linguagens multiparadigma}
\begin{flushright}
  \footnotesize
  \textit{\textbf{1: Paradigmas de programação}}\\
  \textit{c01-conceitos-paradigmas-04}\\

\end{flushright}
Linguagens modernas de amplo uso na indústria, como Python e C++, são frequentemente classificadas como \textbf{multiparadigma}. O que esta classificação indica do ponto de vista da arquitetura de software?

\begin{enumerate}
  \item Indica que a linguagem possui simultaneamente um interpretador em tempo real e um compilador \emph{Ahead-of-Time} (AOT).

  \item Indica que a linguagem exige que o desenvolvedor escolha um único paradigma no momento da criação do projeto, desabilitando permanentemente os comandos dos demais paradigmas.

  \item Indica que a linguagem fornece suporte nativo a técnicas e abstrações de múltiplos paradigmas (como o estruturado, o orientado a objetos e o funcional), permitindo ao desenvolvedor mesclar essas abordagens.

  \item Indica que o compilador da linguagem foi escrito utilizando um paradigma distinto daquele que é exposto aos programadores que a utilizam.

  \item Indica que a linguagem é capaz de compilar o mesmo código-fonte para múltiplas arquiteturas de hardware diferentes (como x86, ARM e RISC-V).


\end{enumerate}

\questionTitle{44}{Acumuladores e avaliação preguiçosa}
\begin{flushright}
  \footnotesize
  \textit{\textbf{9: Funções recursivas}}\\
  \textit{c09-eficiencia-thunks-01}\\

\end{flushright}
Em Haskell, por que uma função com acumulador ainda pode consumir mais memória do que o esperado?

\begin{enumerate}
  \item Porque acumuladores são sempre armazenados como listas infinitas.

  \item Porque o tipo \haskellinline|Integer| impede qualquer otimização do compilador.

  \item Porque Haskell não permite que acumuladores sejam avaliados.

  \item Porque a avaliação preguiçosa pode adiar o cálculo do acumulador, formando expressões suspensas, ou \emph{thunks}.

  \item Porque funções com acumulador não podem ter casos base.


\end{enumerate}

\questionTitle{45}{Precisão arbitrária do Integer}
\begin{flushright}
  \footnotesize
  \textit{\textbf{4: Tipos de Dados}}\\
  \textit{c04-tipos-int-integer-04}\\

\end{flushright}
Se você precisar calcular e armazenar o fatorial de 1000 ($1000!$, um número com mais de 2500 dígitos), qual tipo deve usar e por quê?

\begin{enumerate}
  \item \haskellinline{Integer}, pois é um tipo de precisão arbitrária que pode representar números inteiros de qualquer tamanho.

  \item \haskellinline{Double}, pois a notação científica é a única forma de processar números tão grandes.

  \item \haskellinline{Integer}, pois é o único tipo que aceita operações recursivas profundas.

  \item Ambos falharão; é necessário usar uma biblioteca externa como \haskellinline{BigMath}.

  \item \haskellinline{Int}, pois em Haskell o tipo \haskellinline{Int} escala automaticamente para evitar \emph{overflow}.


\end{enumerate}

\questionTitle{46}{Composição de take, drop e sum}
\begin{flushright}
  \footnotesize
  \textit{\textbf{6: Estruturas de dados básicas}}\\
  \textit{c06-listas-operacoes-01}\\

\end{flushright}
Analise a expressão abaixo.
\begin{minted}{haskell}
let xs = [4..16]
    ys = take 6 (drop 1 xs)
in sum ys
\end{minted}
Qual é o valor final calculado?

\begin{enumerate}
  \item \haskellinline{ 45 }

  \item \haskellinline{ 39 }

  \item A expressão resulta em erro.

  \item \haskellinline{ 33 }

  \item \haskellinline{ 130 }


\end{enumerate}

\questionTitle{47}{Forma fundamental de uma lista}
\begin{flushright}
  \footnotesize
  \textit{\textbf{6: Estruturas de dados básicas}}\\
  \textit{c06-listas-estrutura-01}\\

\end{flushright}
Qual é a forma fundamental da lista \haskellinline{['a', 'b', 'c']} usando apenas o construtor \haskellinline{(:)} e a lista vazia?

\begin{enumerate}
  \item \haskellinline{'a' : 'b' : 'c'}

  \item \haskellinline{[] : 'a' : 'b' : 'c'}

  \item \haskellinline{['a'] : ['b'] : ['c']}

  \item \haskellinline{('a', 'b', 'c')}

  \item \haskellinline{'a' : ('b' : ('c' : []))}


\end{enumerate}

\questionTitle{48}{Status calculado localmente}
\begin{flushright}
  \footnotesize
  \textit{\textbf{5: Expressão condicional}}\\
  \textit{c05-where-avaliacao-01}\\

\end{flushright}
Analise a função:
\begin{haskellcode}
processa x
  | status == "par"   = x * 2
  | status == "impar" = x + 1
  where
    status = if even x then "par" else "impar"
\end{haskellcode}
Qual é o resultado de \haskellinline|processa 7|?

\begin{enumerate}
  \item \haskellinline|8|

  \item \haskellinline|7|

  \item Ocorre um erro, pois \haskellinline|status| só estaria visível depois da cláusula \haskellinline|where|.

  \item Ocorre um erro de tipo na comparação com \haskellinline|"impar"|.

  \item \haskellinline|14|


\end{enumerate}

\questionTitle{49}{Aplicações de missão crítica (Legado Refatorado)}
\begin{flushright}
  \footnotesize
  \textit{\textbf{1: Paradigmas de programação}}\\
  \textit{c01-engenharia-software-01}\\

\end{flushright}
Você foi encarregado de arquitetar o \emph{backend} de um sistema bancário de larga escala. O sistema processa milhões de transações diárias que passam por validações, conversões e cálculos de risco. A \textbf{correção matemática} e a \textbf{auditabilidade} do fluxo de dados são características críticas inegociáveis.

Entre as opções, qual paradigma de programação fornece as abstrações arquiteturais mais seguras para este domínio específico e por quê?

\begin{enumerate}
  \item O paradigma funcional, porque a imutabilidade estrita e a ausência de efeitos colaterais tornam o fluxo de dados totalmente explícito e previsível, facilitando a auditoria e as provas lógicas.

  \item O paradigma imperativo procedural, porque a manipulação direta e mutável do estado da memória é a única forma técnica de garantir a performance de I/O exigida por bancos de dados modernos.

  \item A escolha do paradigma é completamente irrelevante para características de software como auditabilidade e correção matemática; essas métricas dependem exclusivamente da linguagem SQL utilizada.

  \item O paradigma puramente orientado a objetos, porque ocultar transações dentro de objetos de estado opaco garante que auditores externos não consigam corromper os dados acidentalmente.

  \item O paradigma lógico, porque todo o ecossistema bancário pode ser facilmente abstraído como um banco de dados Prolog e resolvido instantaneamente via inferência de axiomas.


\end{enumerate}

\questionTitle{50}{Uso de restrição Show para produzir String}
\begin{flushright}
  \footnotesize
  \textit{\textbf{7: Polimorfismo}}\\
  \textit{cap07-adhoc-show-vs-parametrico}\\

\end{flushright}
Por que a função \haskellinline{show} possui uma restrição de classe de tipo?

\begin{minted}{haskell}
show :: Show a => a -> String
\end{minted}

\begin{enumerate}
  \item Porque toda função que retorna \haskellinline{String} precisa necessariamente pertencer à classe \haskellinline{Show}.
  \item Porque \haskellinline{show} é uma função monomórfica com uma implementação separada apenas para \haskellinline{Int}.
  \item Porque \haskellinline{String} é uma subclasse de \haskellinline{Show}.
  \item Porque converter um valor arbitrário para \haskellinline{String} depende de uma implementação específica de \haskellinline{show} para o tipo do valor.
  \item Porque a restrição \haskellinline{Show a} permite converter automaticamente \haskellinline{String} de volta para \haskellinline{a}.

\end{enumerate}

\questionTitle{51}{Papel dos parênteses em aplicação de função}
\begin{flushright}
  \footnotesize
  \textit{\textbf{3: Expressões e definições}}\\
  \textit{c03-precedencia-associatividade-04}\\

\end{flushright}
Em Haskell, parênteses agrupam expressões, mas não fazem parte da sintaxe obrigatória de chamada de função. Qual alternativa é equivalente a \haskellinline|f (g x) y|?

\begin{enumerate}
  \item \haskellinline|f ((g x) y)|.

  \item \haskellinline|(f (g x)) y|.

  \item \haskellinline|f g x y|.

  \item \haskellinline|f (g (x y))|.

  \item \haskellinline|(f g) (x y)|.


\end{enumerate}

\questionTitle{52}{Cláusula else obrigatória}
\begin{flushright}
  \footnotesize
  \textit{\textbf{5: Expressão condicional}}\\
  \textit{c05-if-regras-01}\\

\end{flushright}
Em Haskell, por que a expressão abaixo é incorreta?
\begin{haskellcode}
if x > 0 then x
\end{haskellcode}

\begin{enumerate}
  \item Porque a condição \haskellinline|x > 0| deveria estar entre parênteses.

  \item Porque o ramo \haskellinline|then| deveria conter uma expressão mais completa do que \haskellinline|x|.

  \item Porque expressões condicionais só podem ser usadas dentro do corpo de funções nomeadas.

  \item Porque o ramo \haskellinline|then| só pode conter valores booleanos.

  \item Porque falta o ramo \haskellinline|else|, obrigatório para que a expressão sempre produza um valor.


\end{enumerate}

\questionTitle{53}{Valor opcional}
\begin{flushright}
  \footnotesize
  \textit{\textbf{6: Estruturas de dados básicas}}\\
  \textit{c06-maybe-conceito-01}\\

\end{flushright}
Qual opção representa corretamente a ideia central da estrutura \haskellinline{Maybe}?

\begin{enumerate}
  \item Ela representa tuplas com campos mutáveis.

  \item Ela representa listas de tamanho variável.

  \item Ela é usada apenas para números opcionais.

  \item Ela representa explicitamente a presença ou a ausência de um valor.

  \item Ela substitui o tipo \haskellinline{Bool} em expressões condicionais.


\end{enumerate}

\questionTitle{54}{Hierarquia de Bibliotecas: Hackage vs Prelude}
\begin{flushright}
  \footnotesize
  \textit{\textbf{2: Ambiente de desenvolvimento}}\\
  \textit{c02-ecossistema-ferramentas-04}\\

\end{flushright}
O ecossistema Haskell divide as bibliotecas e módulos em diferentes níveis de acessibilidade. Em relação a essa estrutura, qual é a principal diferença prática entre o \textbf{Prelude} e o repositório \textbf{Hackage}?

\begin{enumerate}
  \item O Prelude contém as bibliotecas fechadas pagas por licença comercial, enquanto o Hackage é o acervo restrito a componentes de código livre (\emph{open-source}).

  \item Não há diferença; ambos são nomes distintos para a mesma ferramenta interna do GHC responsável por otimizar a avaliação preguiçosa.

  \item O Prelude é utilizado exclusivamente para compilar programas gráficos de janela, enquanto o Hackage é restrito ao desenvolvimento de aplicações de linha de comando.

  \item O Prelude precisa ser importado manualmente em todo arquivo \texttt{.hs}, ao passo que todas as bibliotecas do Hackage já vêm instaladas e prontas para uso direto no GHCi.

  \item O Prelude é carregado automaticamente pelo compilador contendo funções básicas, enquanto o Hackage é um repositório online de onde se baixa pacotes de terceiros usando ferramentas como o Cabal.


\end{enumerate}

\questionTitle{55}{Where não é expressão}
\begin{flushright}
  \footnotesize
  \textit{\textbf{3: Expressões e definições}}\\
  \textit{c03-definicoes-where-04}\\

\end{flushright}
Por que a forma abaixo não é uma maneira válida de criar uma definição local em Haskell?

\begin{haskellcode}
resultado = where x = 1 in x + 1
\end{haskellcode}

\begin{enumerate}
  \item Porque nomes definidos por \haskellinline|where| precisam começar com letra maiúscula.

  \item Porque \haskellinline|x + 1| deveria ser escrito obrigatoriamente como \haskellinline|(+ x 1)|.

  \item Porque \haskellinline|where| só pode ser usado dentro do GHCi, nunca em arquivos fonte.

  \item Porque \haskellinline|where| não permite definir valores numéricos.

  \item Porque \haskellinline|where| não é uma expressão do tipo \haskellinline|where ... in ...|; ele é uma cláusula anexada a uma equação.


\end{enumerate}

\questionTitle{56}{Abreviações sintáticas de listas e strings}
\begin{flushright}
  \footnotesize
  \textit{\textbf{3: Expressões e definições}}\\
  \textit{c03-comentarios-valores-04}\\

\end{flushright}
Qual alternativa apresenta uma equivalência correta de açúcar sintático em Haskell?

\begin{enumerate}
  \item \haskellinline|"abc"| é uma abreviação para \haskellinline|['a','b','c']|.

  \item \haskellinline|(1,2,3)| é uma abreviação para a lista \haskellinline|[1,2,3]|.

  \item \haskellinline|[]| é uma abreviação para a tupla vazia \haskellinline|()|.

  \item \haskellinline|[1,2]| é uma abreviação para a tupla \haskellinline|(1,2)|.

  \item \haskellinline|"abc"| é uma abreviação para \haskellinline|["a","b","c"]|.


\end{enumerate}

\questionTitle{57}{Limitações de uma assinatura totalmente paramétrica}
\begin{flushright}
  \footnotesize
  \textit{\textbf{7: Polimorfismo}}\\
  \textit{cap07-param-limitacao-a-para-b}\\

\end{flushright}
Suponha que desejamos implementar uma função total, sem usar \haskellinline{undefined}, \haskellinline{error}, exceções ou não terminação:

\begin{minted}{haskell}
misterio :: a -> b
\end{minted}

Qual conclusão é tecnicamente adequada?

\begin{enumerate}
  \item A função pode comparar o valor recebido com qualquer outro valor de tipo \haskellinline{b} e escolher um resultado.
  \item Não há implementação total e informativa para essa assinatura, pois a função recebe apenas um valor de tipo \haskellinline{a} e precisa produzir um valor de um tipo arbitrário \haskellinline{b}.
  \item A implementação \haskellinline{misterio x = x} é válida, pois toda variável de tipo pode ser convertida automaticamente para qualquer outra.
  \item Basta exigir \haskellinline{Show a}; então o resultado de tipo \haskellinline{b} pode ser obtido com \haskellinline{show}.
  \item A assinatura força \haskellinline{a} e \haskellinline{b} a serem o mesmo tipo, embora tenham nomes diferentes.

\end{enumerate}

\questionTitle{58}{Equações como definições}
\begin{flushright}
  \footnotesize
  \textit{\textbf{3: Expressões e definições}}\\
  \textit{c03-ghci-definicoes-04}\\

\end{flushright}
Analise a definição:

\begin{haskellcode}
quadrado x = x * x
\end{haskellcode}

Qual alternativa descreve corretamente sua estrutura?

\begin{enumerate}
  \item A definição é inválida porque toda função Haskell deve ter seus parâmetros entre parênteses.

  \item \haskellinline|quadrado| é o nome definido, \haskellinline|x| é um parâmetro formal e \haskellinline|x * x| é a expressão do lado direito da equação.

  \item \haskellinline|quadrado x| é uma tupla formada por dois nomes, e \haskellinline|x * x| é um comentário.

  \item \haskellinline|quadrado| é um construtor, pois aparece à esquerda do sinal de igualdade.

  \item \haskellinline|x * x| é avaliado imediatamente para um único número no momento em que o arquivo é carregado.


\end{enumerate}

\questionTitle{59}{Anotação de tipo em read}
\begin{flushright}
  \footnotesize
  \textit{\textbf{8: Programas interativos}}\\
  \textit{c08-conversao-validacao-03}\\

\end{flushright}
Por que é comum escrever uma anotação como \haskellinline|read linha :: Double| ao converter uma entrada textual?

\begin{enumerate}
  \item Porque \haskellinline|read| sempre retorna \haskellinline|Double|, mas a anotação melhora a legibilidade.

  \item Porque a anotação força o programa a tratar exceções de conversão automaticamente.

  \item Porque \haskellinline|read| é polimórfica. A anotação informa ao compilador qual tipo deve ser produzido pela conversão.

  \item Porque sem anotação \haskellinline|read| lê diretamente da entrada padrão.

  \item Porque \haskellinline|read| só funciona se o tipo final for escrito dentro de uma string.


\end{enumerate}

\questionTitle{60}{let não executa getLine}
\begin{flushright}
  \footnotesize
  \textit{\textbf{8: Programas interativos}}\\
  \textit{c08-blocos-do-03}\\

\end{flushright}
Considere:

\begin{minted}{haskell}
saudacao :: IO ()
saudacao = do
  putStrLn "Qual o seu nome?"
  let nome = getLine
  putStrLn "Olá!"
\end{minted}

O que acontece quando \haskellinline|saudacao| é executada?

\begin{enumerate}
  \item O programa imprime a ação \haskellinline|getLine| como texto antes de imprimir \haskellinline|"Olá!"|.

  \item O programa imprime a pergunta, aguarda o nome e imprime \haskellinline|"Olá!"|.

  \item O programa imprime a pergunta e depois imprime \haskellinline|"Olá!"|, mas não aguarda a digitação do nome. A linha com \haskellinline|let| apenas associa \haskellinline|nome| à ação \haskellinline|getLine|.

  \item O programa não compila, pois \haskellinline|let| é sintaticamente proibido dentro de \haskellinline|do|.

  \item O programa entra em laço infinito porque \haskellinline|getLine| fica esperando entrada indefinidamente.


\end{enumerate}

\questionTitle{61}{Responsabilidades no fechamento de notas}
\begin{flushright}
  \footnotesize
  \textit{\textbf{10: Ações de E/S recursivas}}\\
  \textit{c10-separacao-puro-04}\\

\end{flushright}
Em um programa de fechamento de notas, uma boa organização separa ações de E/S e processamento puro. Qual distribuição de responsabilidades é mais adequada?

\begin{enumerate}
  \item Usar ações de E/S para ler dados e exibir resultados, deixando cálculos como média, nota final, situação e percentuais em funções puras.

  \item Usar \haskellinline|return| para converter todas as funções puras em ações \haskellinline|IO|, mesmo quando não há interação externa.

  \item Colocar leitura, cálculo de nota final, classificação e impressão em uma única função pura, sem tipo \haskellinline|IO|.

  \item Evitar listas, pois programas com \haskellinline|IO| não podem armazenar dados lidos para processamento posterior.

  \item Fazer com que cada função pura leia novamente a matrícula e as notas do aluno para evitar passagem de parâmetros.


\end{enumerate}

\questionTitle{62}{Significado do operador ::}
\begin{flushright}
  \footnotesize
  \textit{\textbf{4: Tipos de Dados}}\\
  \textit{c04-assinatura-operador-01}\\

\end{flushright}
Na declaração \haskellinline{minhaFuncao :: Int -> Bool}, o que o operador \haskellinline{::} significa?

\begin{enumerate}
  \item \emph{tem o tipo de}

  \item \emph{é definido como}

  \item \emph{é um subtipo de}

  \item \emph{pertence à classe}

  \item \emph{retorna}


\end{enumerate}

\questionTitle{63}{Guarda sem cobertura completa}
\begin{flushright}
  \footnotesize
  \textit{\textbf{9: Funções recursivas}}\\
  \textit{c09-fundamentos-cobertura-01}\\

\end{flushright}
Considere a definição:
\begin{haskellcode}
h :: Integer -> Integer
h n
  | n == 0 = 1
  | n > 0  = n * h (n - 1)
\end{haskellcode}
Qual alternativa descreve corretamente o comportamento de \haskellinline|h (-2)|?

\begin{enumerate}
  \item O resultado é \haskellinline|-2|, pois a função devolve o argumento quando ele é negativo.

  \item A avaliação entra em recursão infinita, pois a guarda \haskellinline|n > 0| também vale para negativos.

  \item O resultado é \haskellinline|1|, pois a primeira guarda é usada para qualquer valor não positivo.

  \item Ocorre erro de tipo, pois guardas booleanas não podem ser usadas com \haskellinline|Integer|.

  \item Nenhuma guarda é satisfeita para \haskellinline|(-2)|; portanto, a avaliação falha por ausência de caso aplicável.


\end{enumerate}

\questionTitle{64}{Definição de tipo em Haskell}
\begin{flushright}
  \footnotesize
  \textit{\textbf{4: Tipos de Dados}}\\
  \textit{c04-conceito-definicao-01}\\

\end{flushright}
Em Haskell, um \textbf{tipo} pode ser melhor definido como:

\begin{enumerate}
  \item Uma coleção de valores que compartilham propriedades e operações comuns.

  \item Uma palavra-chave reservada para declaração de constantes.

  \item Uma função que recebe argumentos e retorna um resultado.

  \item Um nome dado a uma variável que armazena um valor numérico.

  \item Um comentário especial que descreve o propósito de uma função.


\end{enumerate}

\questionTitle{65}{String vs IO String}
\begin{flushright}
  \footnotesize
  \textit{\textbf{8: Programas interativos}}\\
  \textit{c08-tipos-acoes-03}\\

\end{flushright}
Considere o trecho:

\begin{minted}{haskell}
nome = getLine
\end{minted}

Sem nenhuma outra informação de contexto, qual é o tipo de \haskellinline|nome|?

\begin{enumerate}
  \item \haskellinline|IO ()|, pois toda ação de E/S sempre produz o valor unitário.

  \item \haskellinline|String -> IO String|, pois \haskellinline|getLine| precisa receber o prompt como argumento.

  \item \haskellinline|String|, pois \haskellinline|getLine| já executa a leitura no momento da definição.

  \item \haskellinline|Char|, pois a entrada padrão lê apenas um caractere por vez.

  \item \haskellinline|IO String|, pois \haskellinline|getLine| é uma ação de E/S que, quando executada, produz uma \haskellinline|String|.


\end{enumerate}

\questionTitle{66}{Consequência da falta de cobertura}
\begin{flushright}
  \footnotesize
  \textit{\textbf{5: Expressão condicional}}\\
  \textit{c05-guardas-exaustividade-01}\\

\end{flushright}
O que acontece se uma função com guardas é chamada com uma entrada para a qual nenhuma guarda resulta em \haskellinline|True|, e não existe outra equação aplicável?

\begin{enumerate}
  \item A função retorna o valor da última expressão escrita na definição.

  \item O compilador sempre rejeita o programa com erro de compilação.

  \item A função passa a retornar automaticamente \haskellinline|undefined|.

  \item Ocorre um erro em tempo de execução, porque a definição não cobre aquela entrada.

  \item O avaliador repete a verificação das guardas até alguma se tornar verdadeira.


\end{enumerate}

\questionTitle{67}{Erro de tipo ao retornar uma ação}
\begin{flushright}
  \footnotesize
  \textit{\textbf{10: Ações de E/S recursivas}}\\
  \textit{c10-return-tipos-03}\\

\end{flushright}
Considere uma tentativa de escrever uma contagem regressiva recursiva:

\begin{haskellcode}
contagem :: Int -> IO ()
contagem n =
  if n <= 0
    then return ()
    else do
      print n
      return (contagem (n - 1))
\end{haskellcode}

Qual é o problema principal dessa definição?

\begin{enumerate}
  \item No ramo recursivo, \haskellinline|return (contagem (n - 1))| cria uma ação que produz outra ação, levando a um tipo como \haskellinline|IO (IO ())| em vez de executar a chamada recursiva.

  \item A função \haskellinline|return| imprime a ação \haskellinline|contagem (n - 1)| em vez de chamá-la.

  \item A condição \haskellinline|n <= 0| deveria ser substituída por \haskellinline|n == 0|, pois guardas não aceitam desigualdades.

  \item O problema é que ações de tipo \haskellinline|IO ()| nunca podem ser definidas recursivamente.

  \item A chamada \haskellinline|print n| precisa ficar depois da chamada recursiva para que o programa compile.


\end{enumerate}

\questionTitle{68}{Avaliação direta de if}
\begin{flushright}
  \footnotesize
  \textit{\textbf{5: Expressão condicional}}\\
  \textit{c05-if-avaliacao-01}\\

\end{flushright}
Qual é o valor da expressão abaixo?
\begin{haskellcode}
if even 6 then 10 else 20
\end{haskellcode}

\begin{enumerate}
  \item \haskellinline|"10"|

  \item Ocorre um erro de sintaxe.

  \item \haskellinline|10|

  \item \haskellinline|True|

  \item \haskellinline|20|


\end{enumerate}

\questionTitle{69}{Quando guardas são mais idiomáticas}
\begin{flushright}
  \footnotesize
  \textit{\textbf{5: Expressão condicional}}\\
  \textit{c05-if-vs-guardas-01}\\

\end{flushright}
Em qual situação o uso de guardas tende a ser mais idiomático e legível do que \haskellinline|if-then-else| aninhado?

\begin{enumerate}
  \item Quando se deseja evitar qualquer comparação booleana.

  \item Quando os ramos do \haskellinline|if| têm tipos diferentes.

  \item Quando a função não recebe argumentos.

  \item Quando a função trata vários casos sequenciais, cada um coberto por uma guarda.

  \item Quando há apenas um teste booleano simples, com dois resultados possíveis.


\end{enumerate}


\end{multicols}
\makeanswersheet{UFOP}{BCC222: Programação Funcional}{José Romildo}{2026/1}{Primeira Prova}{04}{69}{25}{7}
\gabarito{04}{A,C,B,D,B,B,C,E,A,C,C,A,A,E,E,E,D,C,C,E,C,C,A,C,A,A,A,D,C,C,C,A,E,B,E,A,E,E,C,E,E,B,C,D,A,B,E,A,A,D,B,E,D,E,E,A,B,B,C,C,A,A,E,A,E,D,A,C,D}
\end{document}
